\chapter{Conclusion}
\label{chap:conclusion}

\section{Summary}
\label{sec:conclusion_summary}

This thesis presented TF-Rehancer, a time-frequency-domain speech enhancement architecture for computationally controlled full-band processing. The main idea is low-band-guided full-band refinement: TF-Rehancer concentrates heavier modeling on low-band speech structure, keeps observed upper-band evidence in the decoder path, updates high-band decoder tokens using low-band encoder context, and estimates the enhanced spectrum through complex time-frequency filtering over the original noisy STFT.

The 48 kHz VoiceBank+DEMAND results show that the DNS4-trained online TF-Rehancer configuration obtains competitive full-reference scores relative to published PercepNet and DeepFilterNet2 references. The 16 kHz VoiceBank+DEMAND benchmark shows a competitive quality-efficiency trade-off. The compact FastGRU-based setting uses fewer model-side MACs than the paper-reported FastEnhancer-M reference at a similar parameter scale. It reports higher PESQ, SI-SDR, and DNSMOS P.835 values, while FastEnhancer-M remains slightly stronger in SCOREQ and ESTOI. In the controlled 48 kHz ablation setting, direct spectrum mapping degrades full-reference, spectral, and non-intrusive metrics compared with complex TF filtering, suggesting that filtering-based output estimation is more reliable for the proposed architecture. The decoder MHCA ablation shows that decoder cross-attention is more effective than removing it for the main full-reference and spectral trends. The low-band cutoff analysis further shows that 5 kHz provides a balanced operating point between computational cost, waveform fidelity, and full-band spectral fidelity.

Overall, these results support low-band-guided full-band refinement as a computationally controlled design direction for full-band speech enhancement. The architecture preserves observed upper-band information, avoids heavy full-band encoding, and uses full-resolution complex TF filtering to produce the enhanced spectrum.

\section{Limitations and Future Work}
\label{sec:conclusion_future}

The 48 kHz comparison includes paper-reported full-band references. TF-Rehancer is evaluated locally. A fully aligned protocol with locally re-evaluated baselines remains necessary for stronger external full-band claims. The 48 kHz ablation experiments in this work are designed for controlled component analysis. Future work should therefore evaluate TF-Rehancer and competing full-band systems under a shared training and evaluation protocol, including aligned metric computation and runtime measurement.

Several architectural questions also remain open. The present ablation supports complex TF filtering as the output estimator and characterizes the effect of decoder MHCA. Further isolating the lightweight high-band adaptation from the subsequent decoder refinement would clarify this component, for example by removing or replacing the local high-band adapter while preserving the observed high-band embedding. In addition, the low-band cutoff was selected as a balanced operating point under the 48 kHz ablation setting; broader datasets, noise conditions, and sampling rates may require adaptive or data-dependent cutoff strategies.

Finally, practical deployment requires further validation beyond model-side MACs. Future work should include real-time runtime profiling, latency analysis, and evaluation on device-oriented scenarios. Extension to multi-channel or reverberant conditions is another possible direction.

These extensions would further clarify the role of low-band-guided full-band refinement in practical speech enhancement.
